\chapter{Presupuesto de materiales}
\label{anx:presupuesto}

La Tabla~\ref{tab:presupuesto} recoge el coste de los componentes del
instrumento descrito en el capítulo~\ref{cap:hardware}. Las cantidades son las
que hay montadas en el aparato terminado, no las compradas, que en algunos casos
son mayores porque se adquirieron unidades de repuesto.

\begin{table}[htbp]
    \centering
    \caption{Presupuesto de materiales del instrumento.}
    \label{tab:presupuesto}
    \begin{tabular}{llrrr}
        \toprule
        \textbf{Referencia} & \textbf{Componente} & \textbf{Cant.} &
        \textbf{€/ud} & \textbf{Importe} \\
        \midrule
        Daisy Seed 1.2      & Placa de desarrollo de audio        & 1 & 53,25 & 53,25 \\
        ESP32-2432S028      & Placa ESP32 con pantalla táctil     & 1 & 12,51 & 12,51 \\
        ESP32-S3-N16R8      & Placa de desarrollo ESP32-S3        & 1 &  4,74 &  4,74 \\
        MIDI Shield         & Shield MIDI para Arduino            & 1 &  4,73 &  4,73 \\
        Veroboard           & Placa de tiras 100$\times$160~mm    & 1 &  3,36 &  3,36 \\
        POT10k-BLin         & Potenciómetro lineal 10~k$\Omega$   & 6 &  0,29 &  1,71 \\
        OLED 0,96in         & Pantalla OLED de 0,96 pulgadas      & 1 &  1,31 &  1,31 \\
        Switch ON-OFF-ON    & Conmutador de palanca               & 2 &  0,27 &  0,54 \\
        Jack 3,5~mm         & Jack de audio hembra mono           & 1 &  0,33 &  0,33 \\
        Tira 22 pines       & Tira de pines hembra de 2,54~mm     & 2 &  0,15 &  0,30 \\
        Tira 20 pines       & Tira de pines hembra de 2,54~mm     & 2 &  0,15 &  0,30 \\
        R3,3                & Resistencia fusible 3,3~$\Omega$ 1~W & 1 & 0,28 &  0,28 \\
        C470u-ELEC          & Condensador electrolítico 470~µF    & 1 &  0,24 &  0,24 \\
        9V Socket           & Conector de alimentación de 2,1~mm  & 1 &  0,16 &  0,16 \\
        C100n-CER           & Condensador cerámico 100~nF         & 8 &  0,005 & 0,04 \\
        1N5817              & Diodo Schottky                      & 1 &  0,013 & 0,01 \\
        R10k                & Resistencia de 10~k$\Omega$ 1/4~W   & 2 &  0,004 & 0,01 \\
        \midrule
        \multicolumn{4}{l}{\textbf{Coste de materiales}} & \textbf{83,82} \\
        \bottomrule
    \end{tabular}
\end{table}

Tres precisiones sobre cómo está construida la tabla. La primera es que los
importes no incluyen mermas, es decir, el material desechado durante el montaje;
por eso la pantalla OLED figura con una unidad y no con las dos que se
compraron, ya que la primera se rompió al montarla. La segunda es que se han
omitido los componentes cuyo importe, una vez descontada la merma, redondea a
cero céntimos. La tercera es que el presupuesto recoge los componentes
electrónicos y la placa que los soporta, y quedan fuera el alimentador externo
de 230~V a 5~V, el filamento con el que se imprimieron las tres piezas de la
carcasa, el cable del mazo de panel y la tornillería, de modo que el coste real
del conjunto es algo mayor que el que aquí figura.

Conviene señalar que dos de las diecisiete líneas concentran el 78\,\% del
importe, y son las dos placas de desarrollo con procesador de audio y pantalla
táctil. Todo lo demás —los seis potenciómetros, las dos pantallas, los
conectores, la placa de tiras y la electrónica pasiva completa— suma menos de
veinte euros. Es un resultado coherente con el condicionante de partida del
apartado~\ref{sec:requisitos}, que exigía construir el instrumento sobre
hardware comercial de bajo coste sin fabricar ningún circuito a medida.
